% Pursuing Stacks (cotes 134-2 to 134-8): archive cote 134-7, site page 25 (PDF page 26), typescript page 467.
% typescript. Draft compared: ps6.tex:1160-1224.
% AI-transcribed by Claude Opus 5.5 (claude-opus-5-5, adaptive thinking, effort high), Message Batches API, batch ps3 (msgbatch_01XdFmuVVBAcgc7EFxdtCqsj), 2026-11-12; not yet checked by eye.
\scanpage{134-7/25}
\mspage{467}

The latter has \xout{a canonical augmentation} $\pi_i\otimes_k k'$ as homology module in
degree i (and zero homology in degree $j<i$), but the homology modules
\xout{\illegible{1}} \xout{\illegible{1}}
in degree $j>i$ need not be zero. In other words, there is a canonical
\uline{augmentation}
B'') \hfill \xout{\unsure{X'xM(i)'x}} $M(i)'_{*}\longrightarrow K(\pi_i\otimes_k k',i)$ \hfill ,
(where the second member is the semisimplicial \xout{object} k'-module defined
by a i-shifted projective resolution of $\pi_i(X)\otimes_k k'$), but this augmentation
need not be a quasiisomorphism, unless we make the relevant flatness as-%
sumptions. To sum up, \uline{the dévissage of} \xout{\illegible{1}}
\[ X'_{*} = X_{*}\otimes_k k' \]
\uline{deduced from a Postnikov dévissage of $X_{*}$ is} \uuline{not} \uline{a Postnikov dévissage of}
$X'_{*}$ , unless we assume either k' flat over k, or the k-modules $\pi_{\over{i}{j}}(X_{*})$
flat. This at the same time solves the \uline{puzzle} raised on page 444 , and
points towards a \uline{serious shortcoming of the} \inserted{(usual)} \uline{Postnikov dévissage} - \uline{namely}
\uline{that it is not compatible with ground ring extension}, or, as we would
say in the language of algebraic geometry, that this construction is not
``geometric'' - a harsh thing to say indeed !

At this point the idea comes up that we may define another devissage,
a lot more natural in the spirit of a theory of ``abelianization'' of homo-%
topy types it would seem, and which is ``geometric'' namely compatible
with ring extension. Here we'll have to work though with the ``prohibiti-%
ve'' abelianization functor
\[ U(k)\longrightarrow Ab_k \quad , \qquad X\longmapsto L(X) \quad ( = \Gamma^{\wedge}_k(M) ) \]
(where M is a k-module ``representing'' the object X in U(k)), as \deleted{we want}
we'll need the functorial embedding
\margin{X} C) \hfill $X\longrightarrow W(L(X))$ \hfill
(where
\[ W:\ Ab_k\longrightarrow U(k) \]
denotes the canonical functor from k-modules to unipotent k-bundles);
this is licit anyhow, if we admit that \xout{under} the canonical map
\[ L_{pt}(X_{*})\longrightarrow L(X_{*}) \]
(cf. page 454 (w)) is a quasiisomorphism i.e. a weak equivalence, for
the complexes $X_{*}$ we are working with. (This of course should hold over
an arbitrary ground ring, not just $\underline{Z}$.) Applying (C) componentwise, we
get for a complex of bundles $X_{*}$ a canonical map into \deleted{the} \inserted{it's} \xout{abelianized}

\margin{25}
